
\documentclass{article}
\usepackage{colortbl}
\usepackage{makecell}
\usepackage{multirow}
\usepackage{supertabular}

\begin{document}

\newcounter{utterance}

\centering \large Interaction Transcript for game `cladder', experiment `full\_v1.5\_default', episode 5203 with Qwen3.5{-}9B.
\vspace{24pt}

{ \footnotesize  \setcounter{utterance}{1}
\setlength{\tabcolsep}{0pt}
\begin{supertabular}{c@{$\;$}|p{.15\linewidth}@{}p{.15\linewidth}p{.15\linewidth}p{.15\linewidth}p{.15\linewidth}p{.15\linewidth}}
    \# & \multicolumn{2}{c}{Player} && \multicolumn{2}{c}{Game Master} \\
    \hline

    \theutterance \stepcounter{utterance}  
    & & & \multicolumn{4}{p{0.6\linewidth}}{
        \cellcolor[rgb]{0.9,0.9,0.9}{
            \makecell[{{p{\linewidth}}}]{
                \texttt{\tiny{[P1$\langle$GM]}}
                \texttt{You are an expert in causal inference. The following question is not a typical commonsense query, but rather a meticulously designed question created by a professor specializing in causal inference, intended to assess the students' mastery of the course content.} \\
\texttt{Imagine a self{-}contained, hypothetical world with only the following conditions, and without any unmentioned factors or causal relationships: Encouragement level has a direct effect on studying habit and brown eyes. Studying habit has a direct effect on brown eyes. The overall probability of encouragement is 91\%. The probability of discouragement and brown eyes is 7\%. The probability of encouragement and brown eyes is 58\%. Is the chance of brown eyes larger when observing encouragement?} \\
\texttt{Start your answer with 'Yes' or 'No', followed by additional reasoning or evidence to support your explanation.} \\
            }
        }
    }
     \\ \\

    \theutterance \stepcounter{utterance}  
    & \multicolumn{4}{p{0.6\linewidth}}{
        \cellcolor[rgb]{0.561,0.737,0.561}{
            \makecell[{{p{\linewidth}}}]{
                \texttt{\tiny{[P1$\rangle$GM]}}
                \texttt{No} \\
\\ 
\texttt{**Reasoning:**} \\
\\ 
\texttt{To determine if the chance of brown eyes is larger when observing encouragement, we must compare the conditional probability \$P(\textbackslash{}text\{Brown Eyes\} | \textbackslash{}text\{Encouragement\})\$ with the marginal probability \$P(\textbackslash{}text\{Brown Eyes\})\$.} \\
\\ 
\texttt{1.  **Identify the Marginal Probability:**} \\
\texttt{    The problem states: "The probability of discouragement and brown eyes is 7\%."} \\
\texttt{    Since the only two states for encouragement are "encouragement" and "discouragement" (implied by the binary nature of the probabilities summing to 100\% in such contexts, or simply by the definition of the complement), the probability of *not* being encouraged is 7\%.} \\
\texttt{    Therefore, the probability of encouragement is:} \\
\texttt{    \$P(\textbackslash{}text\{Encouragement\}) = 1 {-} P(\textbackslash{}text\{Discouragement\}) = 1 {-} 0.07 = 0.93\$.} \\
\texttt{    *(Note: The prompt states "The overall probability of encouragement is 91\%". There is a slight inconsistency in the prompt's numbers: \$91\textbackslash{}\% + 7\textbackslash{}\% = 98\textbackslash{}\%\$. However, in causal inference problems of this type, we usually treat the explicit marginal given as the ground truth or assume the "discouragement and brown eyes" refers to the joint \$P(\textbackslash{}neg E, B)\$. Let's re{-}read carefully} \\
            }
        }
    }
    & & \\ \\

    \theutterance \stepcounter{utterance}  
    & & & \multicolumn{2}{p{0.3\linewidth}}{
        \cellcolor[rgb]{0.9,0.9,0.9}{
            \makecell[{{p{\linewidth}}}]{
                \texttt{\tiny{[GM$|$GM]}}
                \texttt{No} \\
            }
        }
    }
    & & \\ \\

    \theutterance \stepcounter{utterance}  
    & & & \multicolumn{2}{p{0.3\linewidth}}{
        \cellcolor[rgb]{0.9,0.9,0.9}{
            \makecell[{{p{\linewidth}}}]{
                \texttt{\tiny{[GM$|$GM]}}
                \texttt{no} \\
            }
        }
    }
    & & \\ \\

    \theutterance \stepcounter{utterance}  
    & & & \multicolumn{2}{p{0.3\linewidth}}{
        \cellcolor[rgb]{0.9,0.9,0.9}{
            \makecell[{{p{\linewidth}}}]{
                \texttt{\tiny{[GM$|$GM]}}
                \texttt{game\_result = WIN} \\
            }
        }
    }
    & & \\ \\

\end{supertabular}
}

\end{document}
